\SourceLeaf{530}beigegeben war und eine Berichtigung von Artikel 139 der Disqu. Arithm. enthielt, ist dort an der betreffenden Stelle eingefügt. Die Anmerkung auf Seite 91 ist einer handschriftlichen Notiz entnommen. Im Übrigen unterscheidet sich die vorliegende Ausgabe von den früheren nur durch die Berichtigung einiger Druckfehler. Die von mir hinzugefügten Einschaltungen sind durch eckige Klammern [\,] gekennzeichnet.
\par Die \textit{Tafel des quadratischen Charakters der Primzahlen} ist nach dem Muster der Tabula II der Disqu. Arithm. gedruckt, die in Artikel 99 beschrieben und (in Artikel 331) vorzugsweise zur Zerlegung von Zahlen verwendet wird. Die Schriftzüge der Handschrift mit dem Titel ‘\textit{Quadratorum numeris primis divisorum residua lateralia}’ ähneln am meisten denen des zweiten Teils der Tafel zur Umwandlung gemeiner Brüche in Dezimalbrüche. Anstelle der waagerechten Striche, die einen quadratischen Rest anzeigen, enthält sie kleine Kreise; diejenigen in benachbarten waagerechten oder senkrechten Reihen sind jeweils durch Linien verbunden. Bei der Korrektur fielen mir mehrere Fehler auf. Bei einer einmaligen Vergleichung mit \textsc{Jacobi}’s \textit{Canon Arithmeticus} fand ich dann 190 Abweichungen in den angegebenen Charakteren und durch direkte Bestimmung Übereinstimmung mit jenen gedruckten Tafeln. Entsprechend ist die Ausgabe hier berichtigt.
\par Der erste Teil der \textit{Tafel zur Umwandlung gemeiner Brüche in Dezimalbrüche} ist hier ähnlich wie Tabula III der Disqu. Arithm. eingerichtet. Er enthält für die Primzahlen und ihre Potenzen $p^\pi$ zwischen 3 und 463 die Mantissen $(1)$, $(2)\dotstwo(0)$ der Dezimalbrüche von $\frac{10\pdot r}{p^\pi}$, $\frac{10\pdot rr}{p^\pi}\dotstwo\frac{10}{p^\pi}$. Dabei bedeutet $r$ die Einheit, sodass $(1)=(2)=\dotstwo(0)$ gilt, wenn 10 eine Primitivwurzel von $p^\pi$ ist; andernfalls bezeichnet $r$ die kleinste derjenigen Primitivwurzeln von $p^\pi$, für die als Basis der Index von 10 seinen kleinsten Wert annimmt. Die von 1 verschiedenen Werte von $r$ habe ich zur Erleichterung des Gebrauchs auf Seite 420 der Tafel beigefügt. Die Handschrift, die auch noch nicht die Unterscheidungsziffern der verschiedenen Perioden enthält, ähnelt äußerlich am meisten der Analysis residuorum und scheint zeitlich vor dem hier als zweiten Teil der ganzen Tafel angeordneten Stück zu liegen.
Dieser zweite Teil enthält für die Primzahlen und ihre Potenz $p^\pi$ zwischen 467 und 997 die Mantissen der Dezimalbrüche von $\frac{100}{p^\pi}$. Die Handschrift gibt die Teiler in absteigender Reihenfolge an und schließt mit den Worten: \textit{Explicitus October} 11. 1795. Im Druck habe ich beim Teiler 191 in Periode $(1)$ die 71. Ziffer hinzugefügt und beim Teiler 829 eine Zahl zwischen der 151. und 152. Ziffer weggelassen.
\par Der erste Teil der \textit{Tafel der Häufigkeit der Primzahlen}, die \textsc{Gauss} selbst in einem Brief (Seite 444) erläutert, gibt die Anzahl der Primzahlen in jeder der ersten 1000 Tausendergruppen an und liegt in einer Handschrift von \textsc{Gauss} vor. Im Nachlass fehlen jedoch die im Brief angedeuteten Abzählungen der Hundertergruppen der ersten Million, die eine bestimmte Anzahl von Primzahlen enthalten. Der andere Teil der Tafel, für die zweite und dritte Million, stammt aus einer Handschrift von \textsc{Goldschmidt} allein. Herr \textsc{Meissel} hat durch Abzählung und durch seine Formel die folgenden Berichtigungen zu Seite 436 und 437 gefunden:
